\chapter{Guide for C Programmers}
\label{ch:c-programmers}

This appendix consolidates the key differences between C and Novus for experienced C programmers transitioning to the language. If you've written AmigaOS code in C, this chapter will help you map familiar concepts to their Novus equivalents.

\section{Philosophy Differences}

Novus and C share a systems programming heritage, but differ in philosophy:

\begin{center}
\begin{tabular}{lp{5.5cm}p{5.5cm}}
\textbf{Aspect} & \textbf{C} & \textbf{Novus} \\
\hline
Safety & Trust the programmer & Verify at compile time \\
Defaults & Mutable, public & Immutable, private \\
Errors & Return codes, \texttt{errno} & \texttt{Result<T, E>} types \\
Memory & Manual \texttt{malloc}/\texttt{free} & RAII with \texttt{Drop} trait \\
Null & Implicit in all pointers & Explicit \texttt{*T} vs non-null \texttt{\&T} \\
\end{tabular}
\end{center}

\section{Variable Declarations}

C declares type first, Novus declares intent first:

\begin{lstlisting}[language=C]
// C
int x = 10;              // Mutable by default
const int Y = 20;        // Immutable
static int z = 30;       // File-local
\end{lstlisting}

\begin{lstlisting}
// ... (Novus equivalent - illustrative)
var x = 10               // Mutable (explicit)
let y = 20               // Immutable (explicit)
const Z: i32 = 30        // Compile-time constant
\end{lstlisting}

Key differences:
\begin{itemize}
    \item Type inference: \texttt{var x = 10} infers \texttt{i32}
    \item No semicolons required (optional)
    \item \texttt{let} bindings cannot be reassigned
    \item \texttt{const} must be computable at compile time
\end{itemize}

\section{Type System}

\subsection{Primitive Types}

\begin{center}
\begin{tabular}{ll|ll}
\textbf{C} & \textbf{Novus} & \textbf{C} & \textbf{Novus} \\
\hline
\texttt{char} & \texttt{i8} or \texttt{u8} & \texttt{unsigned char} & \texttt{u8} \\
\texttt{short} & \texttt{i16} & \texttt{unsigned short} & \texttt{u16} \\
\texttt{int} / \texttt{long} & \texttt{i32} & \texttt{unsigned int} & \texttt{u32} \\
\texttt{long long} & \texttt{i64} & \texttt{unsigned long long} & \texttt{u64} \\
\texttt{float} & \texttt{f32} & \texttt{double} & \texttt{f64} \\
\texttt{void*} & \texttt{*u8} & \texttt{NULL} & \texttt{null} \\
\end{tabular}
\end{center}

Novus also provides fixed-point types not available in standard C:
\begin{itemize}
    \item \texttt{fixed16} --- 8.8 fixed-point (useful for Amiga graphics)
    \item \texttt{fixed32} --- 16.16 fixed-point
\end{itemize}

\subsection{Booleans}

\begin{lstlisting}[language=C]
// C: integers are implicitly boolean
if (ptr) { ... }         // Non-zero = true
if (count) { ... }       // Non-zero = true
int flag = 1;            // "true"
\end{lstlisting}

\begin{lstlisting}
// Novus: explicit bool type, no implicit conversion
if ptr != null { ... }   // Must compare explicitly
if count > 0 { ... }     // Must compare explicitly
let flag: bool = true    // Dedicated bool type
\end{lstlisting}

\subsection{Structs}

\begin{lstlisting}[language=C]
// C: typedef dance required
typedef struct {
    int x;
    int y;
} Point;

Point p = {10, 20};
Point *ptr = &p;
int x = ptr->x;          // Arrow for pointer access
\end{lstlisting}

\begin{lstlisting}
struct Point {
    x: i32,
    y: i32,
}
\end{lstlisting}

Usage:

\begin{lstlisting}
// ... (with Point defined above)
let p = Point { x: 10, y: 20 }
let ptr = &p
let x = ptr.x            // Dot works for everything
\end{lstlisting}

Key differences:
\begin{itemize}
    \item No \texttt{typedef} needed---struct name is the type
    \item Field syntax is \texttt{name: type} (reversed from C)
    \item No \texttt{->} operator---dot (\texttt{.}) auto-dereferences
    \item Struct literals use \texttt{Type \{ field: value \}} syntax
\end{itemize}

\section{Pointers and References}

C has one pointer type; Novus distinguishes three:

\begin{center}
\begin{tabular}{lll}
\textbf{Novus} & \textbf{C Equivalent} & \textbf{Properties} \\
\hline
\texttt{\&T} & \texttt{const T*} & Non-null, immutable \\
\texttt{\&var T} & \texttt{T*} & Non-null, mutable \\
\texttt{*T} & \texttt{T*} & Nullable, requires \texttt{unsafe} \\
\end{tabular}
\end{center}

\begin{lstlisting}[language=C]
// C: All pointers can be null
void process(Player *p);           // Might modify, might be null
void print_info(const Player *p);  // Won't modify, might be null
\end{lstlisting}

\begin{lstlisting}
// Novus: Intent is explicit
fn process(p: &var Player)    // Will modify, never null
fn print_info(p: &Player)     // Won't modify, never null
fn maybe_process(p: *Player)  // Might be null (use in unsafe)
\end{lstlisting}

The call site makes mutability visible:

\begin{lstlisting}
process(&var player)   // Clearly shows modification
print_info(&player)    // Clearly shows read-only
\end{lstlisting}

\section{Functions}

\subsection{Declaration Syntax}

\begin{lstlisting}[language=C]
// C: return type first
int add(int a, int b) {
    return a + b;
}

static int helper(void) { return 42; }  // File-local
\end{lstlisting}

\begin{lstlisting}
// Novus: fn keyword, return type last
fn add(a: i32, b: i32) -> i32 {
    return a + b
}

fn helper() -> i32 { return 42 }  // Private by default
pub fn api() -> i32 { return 0 }  // Explicitly public
\end{lstlisting}

\subsection{Methods}

C doesn't have methods; Novus does:

\begin{lstlisting}[language=C]
// C: Pass struct explicitly
int rect_area(const Rectangle *r) {
    return r->width * r->height;
}
void rect_scale(Rectangle *r, int factor) {
    r->width *= factor;
    r->height *= factor;
}

// Usage
int a = rect_area(&rect);
rect_scale(&rect, 2);
\end{lstlisting}

\begin{lstlisting}
// ... (with Rectangle defined elsewhere - illustrative)
impl Rectangle {
    fn area(&self) -> i32 {
        return self.width * self.height
    }
    fn scale(&var self, factor: i32) {
        self.width = self.width * factor
        self.height = self.height * factor
    }
}
\end{lstlisting}

Usage:

\begin{lstlisting}
// ... (with Rectangle defined above)
let a = rect.area()
rect.scale(2)
\end{lstlisting}

\section{Control Flow}

\subsection{If Statements}

\begin{lstlisting}[language=C]
// C: Parentheses required, braces optional
if (x > 0)
    return 1;
else if (x < 0)
    return -1;
else
    return 0;
\end{lstlisting}

\begin{lstlisting}
// Novus: No parentheses, braces required
if x > 0 {
    return 1
} else if x < 0 {
    return -1
} else {
    return 0
}
\end{lstlisting}

\subsection{Switch vs Match}

\begin{lstlisting}[language=C]
// C: switch with fall-through, break required
switch (cmd) {
    case CMD_READ:
        handle_read();
        break;           // Easy to forget!
    case CMD_WRITE:
    case CMD_UPDATE:     // Fall-through
        handle_write();
        break;
    default:
        handle_unknown();
}
\end{lstlisting}

\begin{lstlisting}
// Novus: match with no fall-through
match cmd {
    CMD_READ => handle_read(),
    CMD_WRITE | CMD_UPDATE => handle_write(),  // Explicit OR
    _ => handle_unknown(),
}
\end{lstlisting}

Key differences:
\begin{itemize}
    \item No fall-through---each arm is independent
    \item No \texttt{break} needed
    \item Use \texttt{|} for multiple patterns
    \item \texttt{\_} is the wildcard (like \texttt{default})
    \item \texttt{match} can be an expression returning a value
\end{itemize}

\subsection{Loops}

\begin{lstlisting}[language=C]
// C: Traditional for loop
for (int i = 0; i < 10; i++) {
    process(i);
}
\end{lstlisting}

\begin{lstlisting}
// Novus: Range-based for (preferred)
for i in 0..10 {
    process(i)
}

// C-style also supported
for var i = 0; i < 10; i++ {
    process(i)
}
\end{lstlisting}

\section{Memory Management}

\subsection{The goto cleanup Pattern}

\begin{lstlisting}[language=C]
// C: Manual cleanup with goto
int do_work(void) {
    void *a = NULL, *b = NULL;
    int result = -1;

    a = AllocMem(100, MEMF_ANY);
    if (!a) goto cleanup;

    b = AllocMem(200, MEMF_ANY);
    if (!b) goto cleanup;

    // ... actual work ...
    result = 0;

cleanup:
    if (b) FreeMem(b, 200);  // Must remember size!
    if (a) FreeMem(a, 100);
    return result;
}
\end{lstlisting}

\begin{lstlisting}
// Novus: RAII handles cleanup automatically
fn do_work() -> Result<(), ExecError> {
    var a = MemoryBlock::try_alloc(100, MEMF_ANY)?
    var b = MemoryBlock::try_alloc(200, MEMF_ANY)?

    // ... actual work ...

    return Result::Ok(())
}  // Both freed automatically, correct order, size tracked
\end{lstlisting}

\subsection{RAII vs Manual Management}

\begin{center}
\begin{tabular}{lp{5cm}p{5cm}}
\textbf{Aspect} & \textbf{C} & \textbf{Novus} \\
\hline
Allocation & \texttt{AllocMem(size, flags)} & \texttt{MemoryBlock::alloc(size, flags)} \\
Deallocation & \texttt{FreeMem(ptr, size)} & Automatic via \texttt{Drop} \\
Size tracking & Manual (error-prone) & Automatic \\
Cleanup order & Manual (LIFO required) & Automatic (guaranteed LIFO) \\
Early return & Requires \texttt{goto} & Just \texttt{return} or \texttt{?} \\
\end{tabular}
\end{center}

\section{Error Handling}

\begin{lstlisting}[language=C]
// C: Return codes and errno
BPTR file = Open("data.txt", MODE_OLDFILE);
if (!file) {
    LONG err = IoErr();
    printf("Error: %ld\n", err);
    return -1;
}
// ... use file ...
Close(file);
\end{lstlisting}

\begin{lstlisting}
// Novus: Result type with ? operator
fn read_data() -> Result<String, DosError> {
    let file = open("data.txt", MODE_OLDFILE)?  // Returns on error
    let data = read_all(&file)?
    return Result::Ok(data)
}  // file closed automatically
\end{lstlisting}

The \texttt{?} operator:
\begin{enumerate}
    \item Checks if result is \texttt{Err}
    \item If error, returns immediately with that error
    \item If success, unwraps the value and continues
\end{enumerate}

\section{Headers and Modules}

\begin{lstlisting}[language=C]
// C: Textual inclusion
#ifndef MYHEADER_H
#define MYHEADER_H

#include <exec/types.h>
#include <proto/exec.h>  // Order can matter!

extern int global_count;
void my_function(void);

#endif
\end{lstlisting}

\begin{lstlisting}
// Novus: Semantic imports
from std::ffi::exec import AllocMem, FreeMem
from std::ffi::amiga_consts import MEMF_FAST

pub fn my_function() {
    // ...
}
\end{lstlisting}

Key differences:
\begin{itemize}
    \item No include guards needed
    \item Import order doesn't matter
    \item Only imported names are in scope
    \item No separate header files
    \item \texttt{pub} explicitly exports (vs C's implicit visibility)
\end{itemize}

\section{Visibility}

\begin{center}
\begin{tabular}{ll}
\textbf{C} & \textbf{Novus} \\
\hline
\texttt{static} (file-local) & (default, no modifier) \\
Non-static + header & \texttt{pub} \\
\texttt{extern} declaration & \texttt{extern} (for FFI) \\
\end{tabular}
\end{center}

Note: Novus inverts C's default. In C, symbols are public unless marked \texttt{static}. In Novus, symbols are private unless marked \texttt{pub}.

\section{Common Intrinsics}

\begin{center}
\begin{tabular}{ll}
\textbf{C} & \textbf{Novus} \\
\hline
\texttt{sizeof(T)} & \texttt{@sizeof(T)} \\
\texttt{\_Alignof(T)} & \texttt{@alignof(T)} \\
\texttt{memset(\&x, 0, sizeof(x))} & \texttt{@zeroed(T)} \\
\texttt{typeof(expr)} (GCC) & \texttt{@typeof(expr)} \\
\end{tabular}
\end{center}

\section{AmigaOS NDK Access}

\textbf{The entire AmigaOS NDK is available from Novus.} Every function, structure, and constant from the 3.9 NDK has been imported and can be called directly. You don't need wrappers---you can use the raw API just like you would in C.

\subsection{FFI Modules}

The NDK bindings are organized by library in \texttt{std::ffi}:

\begin{lstlisting}
// ... (FFI import examples - illustrative)
from std::ffi::exec import AllocMem, FreeMem, FindTask, Wait, Signal
from std::ffi::dos import Open, Close, Read, Write, IoErr, Lock, UnLock
from std::ffi::graphics import BltBitMap, SetAPen, RectFill, Text
from std::ffi::intuition import OpenWindow, CloseWindow, RefreshGList
from std::ffi::gadtools import CreateMenus, FreeMenus, LayoutMenus
from std::ffi::layers import CreateUpfrontLayer, DeleteLayer
from std::ffi::utility import AllocateTagItems, FreeTagItems
\end{lstlisting}

Structures and constants:

\begin{lstlisting}
// ... (structure and constant imports - illustrative)
from std::ffi::amiga_structs import Window, Screen, BitMap
from std::ffi::amiga_consts import MEMF_CHIP, MEMF_FAST, MEMF_CLEAR
\end{lstlisting}

\subsection{Calling Raw NDK Functions}

Raw AmigaOS calls require \texttt{unsafe} blocks because they bypass Novus's safety guarantees:

\begin{lstlisting}
from std::ffi::exec import AllocMem, FreeMem
from std::ffi::amiga_consts import MEMF_CHIP, MEMF_CLEAR

fn allocate_chip_buffer(size: u32) -> *u8 {
    unsafe {
        return AllocMem(size, MEMF_CHIP | MEMF_CLEAR)
    }
}

fn free_chip_buffer(ptr: *u8, size: u32) {
    if ptr != null {
        unsafe {
            FreeMem(ptr, size)
        }
    }
}
\end{lstlisting}

This is identical to what you'd write in C, just with explicit \texttt{unsafe} markers.

\subsection{When to Use Raw vs Wrappers}

\begin{center}
\begin{tabular}{lp{5cm}p{5cm}}
\textbf{Use} & \textbf{Raw NDK} & \textbf{Stdlib Wrappers} \\
\hline
When & Porting C code, one-off calls, no wrapper exists & New code, repeated patterns \\
Safety & Manual (like C) & Automatic (RAII, Result) \\
Cleanup & Manual \texttt{Free*} calls & Automatic via \texttt{Drop} \\
Errors & Check return values & \texttt{Result<T, E>} \\
\end{tabular}
\end{center}

\textbf{Both approaches are valid.} The wrappers exist for convenience and safety, but they're built on the same raw FFI you have access to. If a wrapper doesn't exist for what you need, use the raw API directly.

\subsection{Example: Raw Intuition Window}

Here's opening a window the ``C way'' in Novus:

\begin{lstlisting}
from std::ffi::intuition import OpenWindowTags, CloseWindow
from std::ffi::amiga_consts import WA_Width, WA_Height, WA_Title,
    WA_IDCMP, WA_Flags, WFLG_CLOSEGADGET, WFLG_DRAGBAR,
    IDCMP_CLOSEWINDOW, TAG_DONE

fn main() -> i32 {
    var window: *Window = null

    unsafe {
        window = OpenWindowTags(null,
            WA_Width, 640,
            WA_Height, 480,
            WA_Title, "Raw NDK Window",
            WA_Flags, WFLG_CLOSEGADGET | WFLG_DRAGBAR,
            WA_IDCMP, IDCMP_CLOSEWINDOW,
            TAG_DONE)
    }

    if window == null {
        return 1
    }

    // ... use window ...

    unsafe {
        CloseWindow(window)
    }

    return 0
}
\end{lstlisting}

This is nearly identical to C code---the only differences are:
\begin{itemize}
    \item \texttt{unsafe \{\}} blocks around library calls
    \item \texttt{null} instead of \texttt{NULL}
    \item Novus variable declaration syntax
\end{itemize}

\subsection{BCPL Pointers (DOS Library)}

DOS library functions use BCPL pointers (\texttt{BPTR}/\texttt{BSTR}), which are \texttt{i32} values, not actual pointers:

\begin{lstlisting}
from std::ffi::dos import Open, Close, Read, Write
from std::ffi::amiga_consts import MODE_OLDFILE

fn read_file_raw(path: *u8) -> i32 {
    // DOS file handle is i32 (BPTR), not *File
    var file: i32 = 0

    unsafe {
        file = Open(path, MODE_OLDFILE)
    }

    if file == 0 {
        return -1
    }

    var buffer = [0u8; 1024]
    var bytes_read: i32 = 0

    unsafe {
        bytes_read = Read(file, &buffer as *u8, 1024)
        Close(file)
    }

    return bytes_read
}
\end{lstlisting}

\subsection{Comparison: Raw vs Wrapper}

For comparison, here's the same operation using the stdlib wrapper:

\begin{lstlisting}[language=C]
// C: Direct library calls
struct Window *win = OpenWindowTags(NULL,
    WA_Width, 640,
    WA_Height, 480,
    WA_Title, "My Window",
    TAG_DONE);
if (win) {
    // ... use window ...
    CloseWindow(win);
}
\end{lstlisting}

\begin{lstlisting}
// Novus: RAII wrappers preferred
match WindowBuilder::new()
    .size(640, 480)
    .title("My Window")
    .build()
{
    Result::Ok(window) => {
        // ... use window ...
    },  // Closed automatically
    Result::Err(e) => {
        println(e.message())
    }
}
\end{lstlisting}

\subsection{Hardware Access}

Both languages support direct hardware access, but Novus requires \texttt{unsafe}:

\begin{lstlisting}[language=C]
// C: Direct hardware poke
volatile UWORD *custom = (UWORD *)0xdff000;
custom[0x96/2] = 0x8020;  // DMACON
\end{lstlisting}

\begin{lstlisting}
// ... (Novus unsafe block - illustrative)
unsafe {
    let custom = 0xdff000 as *u16
    *(custom + 0x96 / 2) = 0x8020
}
\end{lstlisting}

\section{Quick Reference Table}

\begin{center}
\begin{tabular}{ll}
\textbf{C} & \textbf{Novus} \\
\hline
\texttt{int x = 10;} & \texttt{var x = 10} \\
\texttt{const int X = 10;} & \texttt{const X: i32 = 10} \\
\texttt{int arr[3] = \{1,2,3\};} & \texttt{let arr = [1, 2, 3]} \\
\texttt{if (ptr)} & \texttt{if ptr != null} \\
\texttt{switch/case/break} & \texttt{match \{ pat => \}} \\
\texttt{for(i=0;i<n;i++)} & \texttt{for i in 0..n} \\
\texttt{ptr->field} & \texttt{ptr.field} \\
\texttt{struct S \{...\};} & \texttt{struct S \{...\}} \\
\texttt{typedef struct} & (not needed) \\
\texttt{malloc/free} & \texttt{MemoryBlock}, \texttt{Vec} \\
\texttt{sizeof(T)} & \texttt{@sizeof(T)} \\
\texttt{NULL} & \texttt{null} \\
\texttt{errno} & \texttt{Result::Err(...)} \\
\texttt{static} & (default) \\
\texttt{\#include} & \texttt{from ... import} \\
\end{tabular}
\end{center}

\section{Migration Tips}

\begin{enumerate}
    \item \textbf{Start with \texttt{var}} --- You can always tighten to \texttt{let} later
    \item \textbf{Use stdlib wrappers} --- \texttt{MemoryBlock}, \texttt{Vec}, \texttt{WindowHandle} instead of raw FFI
    \item \textbf{Embrace \texttt{Result}} --- The \texttt{?} operator makes error handling concise
    \item \textbf{Trust the Drop} --- Let RAII handle cleanup instead of manual \texttt{goto cleanup}
    \item \textbf{Check explicit null} --- \texttt{if ptr != null} instead of \texttt{if (ptr)}
    \item \textbf{Use match} --- More powerful and safer than \texttt{switch}
    \item \textbf{Read the sidebars} --- ``Coming from C'' boxes throughout the guide explain specific differences
\end{enumerate}

\section{What Stays the Same}

Not everything is different! These concepts transfer directly:

\begin{itemize}
    \item Structs are laid out the same way (C ABI compatible)
    \item Pointer arithmetic works the same
    \item Bitwise operations are identical
    \item AmigaOS structures and constants are the same
    \item The calling convention matches (d0/d1/a0/a1 for args)
    \item You can still write inline assembly when needed
\end{itemize}

Novus is designed to feel familiar to C programmers while adding safety. Your AmigaOS knowledge transfers directly---you're just expressing it in a slightly different syntax with more compile-time guarantees.
